\documentclass[11pt,a4paper]{article}

\usepackage[UTF8]{ctex}
\usepackage{geometry}
\geometry{margin=2.2cm}
\usepackage{hyperref}
\usepackage{listings}
\usepackage{xcolor}
\usepackage{enumitem}
\usepackage{booktabs}
\usepackage{longtable}

\hypersetup{
  colorlinks=true,
  linkcolor=blue,
  urlcolor=blue
}

\definecolor{codebg}{RGB}{245,245,245}
\definecolor{codefg}{RGB}{25,25,25}

\lstset{
  backgroundcolor=\color{codebg},
  basicstyle=\ttfamily\small\color{codefg},
  breaklines=true,
  columns=fullflexible,
  frame=single,
  framerule=0.2pt,
  xleftmargin=8pt,
  xrightmargin=8pt,
  aboveskip=10pt,
  belowskip=10pt
}

\title{Windows 上安装与使用 Codex CLI 的新手分步教程\\（含 WSL2、代理配置、项目接入与 skills 工作流）}
\author{currydai}
\date{\today}
\setcounter{secnumdepth}{0}

\begin{document}
\maketitle
% \tableofcontents
\newpage

\section{适用范围与目标}
\subsection{适用范围}
本教程面向 Windows 用户，目标是在本机完成以下能力搭建：
\begin{enumerate}[leftmargin=1.7em]
  \item 安装并运行 Codex CLI；
  \item 在本机项目目录中使用 Codex（读写代码、运行命令、生成文件）；
  \item 与 skills 协同，形成可复用的自动化开发/文档生成工作流；
  \item 解决常见环境问题：WSL 网络与代理、Python 依赖与虚拟环境、权限差异等。
\end{enumerate}

\subsection{为什么推荐 WSL2（官方实践更接近 Linux 开发环境）}
尽管你可以在 Windows 原生终端里使用部分工具，但在实际工程中：
\begin{itemize}[leftmargin=1.5em]
  \item 许多开发工具链（shell 脚本、包管理器、路径惯例、权限模型）在 Linux 上更统一；
  \item skills 往往需要稳定的命令行环境，WSL2 在 Windows 上提供了接近原生的 Linux 体验；
  \item 因此，本教程\textbf{推荐采用 WSL2 + Ubuntu} 作为 Codex CLI 的主要运行环境。
\end{itemize}

\section{准备工作（请先确认这三点）}
\subsection{硬件与系统}
\begin{itemize}[leftmargin=1.5em]
  \item Windows 10/11（建议 Windows 11）；
  \item 建议 8GB 内存及以上；
  \item 可正常联网（如存在网络限制，请参考第 \ref{sec:proxy} 节配置代理）。
\end{itemize}

\subsection{你将会用到的三个终端}
\begin{enumerate}[leftmargin=1.7em]
  \item \textbf{Windows PowerShell（管理员）}：用于启用 WSL2；
  \item \textbf{WSL Ubuntu 终端}：日常开发与运行 Codex CLI；
  \item （可选）\textbf{Windows Terminal}：更好的多标签体验（可同时开 PowerShell 与 Ubuntu）。
\end{enumerate}

\subsection{快速自检命令}
在 WSL Ubuntu 终端中执行：
\begin{lstlisting}
uname -a
pwd
\end{lstlisting}
若能看到 Linux/WSL 信息，且路径如 \texttt{/home/...} 或 \texttt{/mnt/c/...}，说明环境正常。

\section{第 1 步：安装 WSL2 + Ubuntu}
\subsection{启用 WSL2}
以管理员身份打开 PowerShell，执行：
\begin{lstlisting}
wsl --install -d Ubuntu
wsl --update
\end{lstlisting}
如提示重启，请重启 Windows。

\subsection{初始化 Ubuntu}
在开始菜单打开 ``Ubuntu''，首次启动会要求设置：
\begin{itemize}[leftmargin=1.5em]
  \item Linux 用户名（任意）；
  \item Linux 密码（输入不回显属正常）。
\end{itemize}

\section{第 2 步：在 WSL 中安装运行时（Node.js / npm）}
Codex CLI 通常通过 Node 生态安装，因此先安装 Node.js 与 npm。

\subsection{安装 Node.js 与 npm（快速方式）}
在 Ubuntu 终端执行：
\begin{lstlisting}
sudo apt update
sudo apt install -y nodejs npm
node -v
npm -v
\end{lstlisting}

\subsection{（可选）安装更高版本 Node（推荐用于长期开发）}
若你希望使用最新 LTS，可采用 nvm：
\begin{lstlisting}
curl -o- https://raw.githubusercontent.com/nvm-sh/nvm/v0.39.7/install.sh | bash
# 关闭并重新打开终端后：
nvm install --lts
node -v
npm -v
\end{lstlisting}

\section{第 3 步：安装 Codex CLI}
在 Ubuntu 终端执行：
\begin{lstlisting}
npm install -g @openai/codex
codex --version
\end{lstlisting}
若能输出版本号，说明安装成功。

\section{第 4 步：鉴权（登录/密钥配置）}
Codex CLI 通常支持两种方式（以你实际 CLI 提示为准）：
\begin{enumerate}[leftmargin=1.7em]
  \item \textbf{交互式授权}：按 CLI 引导登录/授权；
  \item \textbf{API Key}：设置环境变量供 CLI 使用。
\end{enumerate}

\subsection{方式 A：按 CLI 引导授权}
直接执行：
\begin{lstlisting}
codex
\end{lstlisting}
按照界面提示完成授权即可。

\subsection{方式 B：使用 API Key（通用方式）}
在 WSL 终端中设置环境变量：
\begin{lstlisting}
export OPENAI_API_KEY="sk-xxxxxxxxxxxxxxxx"
\end{lstlisting}
为了永久生效，写入 \texttt{\textasciitilde/.bashrc}：
\begin{lstlisting}
echo 'export OPENAI_API_KEY="sk-xxxxxxxxxxxxxxxx"' >> ~/.bashrc
source ~/.bashrc
\end{lstlisting}

\section{第 5 步：网络与代理配置（若你处于受限网络）}
\label{sec:proxy}
如果出现请求超时、pip/npm 下载失败等情况，优先按本节排查。

\subsection{5.1 验证 WSL 直连是否通}
\begin{lstlisting}
curl -I --max-time 10 https://www.google.com
\end{lstlisting}
若超时，再尝试：
\begin{lstlisting}
curl -I --max-time 10 https://1.1.1.1
\end{lstlisting}

\subsection{5.2 验证代理是否可用（推荐用 -x 显式指定）}
假设本机代理端口为 \texttt{127.0.0.1:51837}：
\begin{lstlisting}
curl -I -x http://127.0.0.1:51837 https://www.google.com
\end{lstlisting}
出现 ``\texttt{200 Connection established}'' 且后续有 HTTP 状态行，说明代理链路正常。

\subsection{5.3 设置全局代理环境变量（建议写入 ~/.bashrc）}
\begin{lstlisting}
export http_proxy="http://127.0.0.1:51837"
export https_proxy="http://127.0.0.1:51837"
export HTTP_PROXY="$http_proxy"
export HTTPS_PROXY="$https_proxy"
\end{lstlisting}
验证：
\begin{lstlisting}
curl -I --max-time 10 https://www.google.com
\end{lstlisting}

\subsection{5.4 npm 与 git 的代理（可选）}
\textbf{npm：}
\begin{lstlisting}
npm config set proxy http://127.0.0.1:51837
npm config set https-proxy http://127.0.0.1:51837
\end{lstlisting}
\textbf{git：}
\begin{lstlisting}
git config --global http.proxy http://127.0.0.1:51837
git config --global https.proxy http://127.0.0.1:51837
\end{lstlisting}

\section{第 6 步：把 Codex CLI 接入你的本机项目（最关键的一步）}
\subsection{6.1 理解路径：Windows 与 WSL 的映射关系}
在 WSL 中，Windows 盘符通常挂载为：
\begin{itemize}[leftmargin=1.5em]
  \item \texttt{C:\textbackslash Users\textbackslash name} $\rightarrow$ \texttt{/mnt/c/Users/name}
  \item \texttt{D:\textbackslash BOOK} $\rightarrow$ \texttt{/mnt/d/BOOK}
\end{itemize}

\subsection{6.2 进入项目目录}
例如项目在 D 盘：
\begin{lstlisting}
cd /mnt/d/BOOK/THE_ROAD_TO_AI
ls
\end{lstlisting}

\subsection{6.3 在项目目录中启动 Codex CLI}
\begin{lstlisting}
cd /mnt/d/BOOK/THE_ROAD_TO_AI
codex
\end{lstlisting}
\textbf{原则：} Codex 的文件读写与命令执行，默认围绕你启动时所在目录展开；在项目根目录启动最清晰。

\subsection{6.4 为什么你会觉得“慢”（官方解释口径）}
常见原因：
\begin{enumerate}[leftmargin=1.7em]
  \item 首次扫描/索引仓库需要时间；
  \item 代理不稳定会导致请求重试；
  \item WSL 访问 \texttt{/mnt/c}、\texttt{/mnt/d} 相比 \texttt{/home} 会更慢。
\end{enumerate}
\textbf{优化建议：} 若追求性能，可将项目复制到 WSL 的 Linux 文件系统（如 \texttt{/home/<user>/projects/...}）再工作。

\section{第 7 步：Codex CLI 的基础用法（新手必读）}
\subsection{7.1 Codex 的工作方式（你需要知道的三件事）}
\begin{enumerate}[leftmargin=1.7em]
  \item 你用自然语言提出任务；
  \item Codex 会提出它计划执行的操作（读文件、改代码、跑命令等）；
  \item 通常需要你确认（approve）后才会实际执行。
\end{enumerate}

\subsection{7.2 推荐的提问格式（更稳定、更可复现）}
建议你在需求中明确写出：
\begin{itemize}[leftmargin=1.5em]
  \item \textbf{目标}：要实现什么；
  \item \textbf{范围}：涉及哪些目录/文件；
  \item \textbf{约束}：不能做什么（例如不能 sudo、不能新增系统依赖）；
  \item \textbf{验收}：运行哪条命令算成功、期望输出是什么。
\end{itemize}

\textbf{示例：}
\begin{lstlisting}
We are in /mnt/d/BOOK/THE_ROAD_TO_AI/report_template.
Goal: generate an A4 PDF report template.
Constraints: no sudo; use only Python venv dependencies.
Acceptance: running "python generate_report.py" outputs report_template.pdf.
\end{lstlisting}

\subsection{7.3 常见操作示例}
\textbf{理解项目：}
\begin{lstlisting}
Explain this repository structure and how to run it.
\end{lstlisting}
\textbf{修复问题：}
\begin{lstlisting}
Find and fix the bug causing a crash when input is empty. Add a minimal test.
\end{lstlisting}
\textbf{新增功能并给验证命令：}
\begin{lstlisting}
Add a --dry-run flag to scripts/build.py, then show how to verify it.
\end{lstlisting}

\subsection{7.4 如何退出 Codex CLI}
通用方式：
\begin{itemize}[leftmargin=1.5em]
  \item 输入 \texttt{/exit} 或 \texttt{/quit}（若支持命令模式）；
  \item 或按 \texttt{Ctrl+C}（中断）/ \texttt{Ctrl+D}（结束输入）退出。
\end{itemize}

\section{第 8 步：skills 是什么？如何在 Codex CLI 中使用 skills？}
\subsection{8.1 skills 的定位（官方表述风格）}
skills 可理解为面向特定任务类别的可复用能力包，用于将常见工作流（如 PDF 生成、自动化测试、部署、截图、表格生成等）标准化、模板化，从而提升可重复性与交付效率。

\subsection{8.2 skills 的典型使用方式}
一般流程：
\begin{enumerate}[leftmargin=1.7em]
  \item 在 Codex 环境中查看可用 skills；
  \item 安装/启用所需 skill；
  \item 在项目目录中发出明确指令：\textbf{使用某个 skill 完成某个任务，并给出验收方式}。
\end{enumerate}

\subsection{8.3 新手推荐的“技能调用提示词模板”}
\begin{lstlisting}
Use the "pdf" skill to generate a clean report template:
- A4 layout, cover page, table of contents
- sections: Summary, Methods, Results, Discussion, Appendix
- header/footer with page number and document title
- output: report_template.pdf
Also provide how to customize the template.
\end{lstlisting}

\section{第 9 步：Python 项目依赖管理（解决你遇到的 externally-managed-environment）}
\subsection{9.1 为什么会报 externally-managed-environment（PEP 668）}
在 Debian/Ubuntu 系统中，系统 Python 由包管理器维护。为了避免系统级 Python 被 pip ``污染''，系统会限制直接向全局环境安装第三方包，从而提示：
\begin{center}
\texttt{error: externally-managed-environment}
\end{center}
\textbf{标准做法：使用虚拟环境 venv。}

\subsection{9.2 一次性正确配置（推荐步骤）}
\textbf{安装 venv 支持（需要 sudo）：}
\begin{lstlisting}
sudo apt update
sudo apt install -y python3-venv python3-pip
\end{lstlisting}

\textbf{在项目目录创建并激活虚拟环境：}
\begin{lstlisting}
cd /mnt/d/BOOK/THE_ROAD_TO_AI/report_template
python3 -m venv .venv
source .venv/bin/activate
python -V
pip -V
\end{lstlisting}

\textbf{在虚拟环境中安装依赖：}
\begin{lstlisting}
pip install -U pip
pip install reportlab
\end{lstlisting}

\subsection{9.3 常见误区（请对照排查）}
\begin{itemize}[leftmargin=1.5em]
  \item \textbf{误区 1：} 在系统环境直接 \texttt{pip install ...}（会触发 PEP 668）。
  \item \textbf{误区 2：} 激活脚本用错：WSL/Linux 用 \texttt{source .venv/bin/activate}，不是 Windows 的 \texttt{.venv\textbackslash Scripts\textbackslash activate}.
  \item \textbf{误区 3：} 在 \texttt{.venv/bin} 目录里执行 \texttt{./activate}（不推荐）。应在项目目录执行 \texttt{source .venv/bin/activate}。
\end{itemize}

\section{第 10 步：从“搭建”到“交付”的完整示例（以 PDF 报告模板为例）}
\subsection{目标}
在项目目录中生成一个标准 A4 PDF 报告模板（封面、目录、章节、页眉页脚），可用于后续自动化报告输出。

\subsection{推荐实施路径（最稳）}
\begin{enumerate}[leftmargin=1.7em]
  \item 进入项目目录；
  \item 创建 venv 并激活；
  \item 安装依赖（reportlab 等）；
  \item 运行脚本生成 PDF；
  \item 用 Codex CLI 固化成可复用模板，并用 skills 进一步标准化。
\end{enumerate}

\subsection{命令清单（可复制执行）}
\begin{lstlisting}
# 1) 进入项目目录
cd /mnt/d/BOOK/THE_ROAD_TO_AI/report_template

# 2) 创建并激活虚拟环境
python3 -m venv .venv
source .venv/bin/activate

# 3) 安装依赖
pip install -U pip
pip install reportlab

# 4) 生成 PDF
python generate_report.py
\end{lstlisting}

\subsection{若出现 Cairo / renderPM / rlPyCairo 相关错误，如何处理}
如果你看到类似：
\begin{itemize}[leftmargin=1.5em]
  \item \texttt{No module named rlPyCairo}
  \item \texttt{No module named \_rl\_renderPM}
  \item \texttt{Dependency lookup for cairo ... failed: pkg-config not found}
\end{itemize}
这通常意味着你启用了 reportlab 的某些渲染后端，需要系统级依赖（pkg-config、cairo、编译工具链等）。

\textbf{两条策略：}
\begin{enumerate}[leftmargin=1.7em]
  \item \textbf{推荐（新手/无 sudo 场景）：} 调整脚本/模板，避免矢量$\rightarrow$PNG 渲染，改用现成 PNG/JPG 占位图；
  \item \textbf{完整（需要 sudo）：} 安装系统依赖（pkg-config/cairo 开发包等）再装 pycairo/rlPyCairo。
\end{enumerate}

\textbf{完整依赖示例（需要 sudo）：}
\begin{lstlisting}
sudo apt install -y pkg-config cmake build-essential libcairo2-dev python3-dev
pip install pycairo rlPyCairo
\end{lstlisting}

\section{第 11 步：建议的日常工作流（最适合新手的固定流程）}
\begin{enumerate}[leftmargin=1.7em]
  \item 打开 WSL Ubuntu；
  \item \textbf{cd 到项目目录}；
  \item \textbf{激活 venv}（Python 项目必做）；
  \item 检查网络/代理（必要时）；
  \item 启动 \texttt{codex}；
  \item 用 ``目标 + 约束 + 验收'' 方式发出任务；
  \item 本地运行验收命令确认结果。
\end{enumerate}

\section{附录 A：Windows 与 WSL 的命令差异速查}
\begin{longtable}{@{}p{0.25\textwidth}p{0.35\textwidth}p{0.35\textwidth}@{}}
\toprule
场景 & Windows (PowerShell) & WSL (Ubuntu Bash) \\
\midrule
进入目录 & \texttt{cd D:\textbackslash BOOK} & \texttt{cd /mnt/d/BOOK} \\
列文件 & \texttt{dir} & \texttt{ls} \\
安装系统包 & \texttt{winget/choco} & \texttt{sudo apt install} \\
激活 venv & \texttt{.venv\textbackslash Scripts\textbackslash activate} & \texttt{source .venv/bin/activate} \\
查看环境变量 & \texttt{\$env:HTTP\_PROXY} & \texttt{echo \$http\_proxy} \\
\bottomrule
\end{longtable}

\section{附录 B：最常见问题（FAQ）}
\subsection*{Q1：我应该先激活 venv 再启动 codex 吗？}
\textbf{建议是：如果项目是 Python 项目，先激活 venv 再启动 codex。}\\
理由：Codex 运行你项目命令时，会继承当前 shell 的环境，包含 Python 路径、依赖、代理等。

\subsection*{Q2：为什么 apt install 报 Permission denied？}
因为 apt 需要 root 权限。请使用：
\begin{lstlisting}
sudo apt install ...
\end{lstlisting}

\subsection*{Q3：如何确认 Codex 已经在我的项目里？}
在启动 Codex 前先确认：
\begin{lstlisting}
pwd
ls
\end{lstlisting}
确保当前目录就是项目根目录或你期望的工作目录。

\end{document}
